% !TeX root = ../main.tex

\begin{abstract}
Este artigo avalia a redução do prazo de um processo de projetos de arquitetura e engenharia por meio da reorganização dos inícios das atividades, mantendo os dados preparados. Na instância completa, com 243 atividades, 1.008 vínculos e 39 equipes, uma formulação inteira resolvida pelo CBC encontrou prazo de 307 dias, mas parou sem provar o ótimo. Por isso, aplicou-se também o esquema serial de geração de cronogramas, combinado ao caminho crítico: o método construtivo encontrou e validou cronogramas de 285 dias quando se impõe que 27 atividades não comecem antes das datas preservadas, e de 187 dias com largadas livres. Em ambas as variantes, os cronogramas atingem o limite inferior de precedências e, portanto, são ótimos para as hipóteses adotadas. A redução de 25 dias (8,1\%) no cenário com liberações é relativa à referência de 310 dias; não constitui previsão validada da duração real. A diferença de 98 dias entre as variantes decorre da imposição conjunta das liberações e não mede diretamente a duração do estudo preliminar ausente. Os resultados comparam planejamentos sob hipóteses fixas.
\end{abstract}

\palavraschave{Sequenciamento de projetos; programação inteira;
restrições de precedência; recursos limitados.}
